% howto.tex - How to Use This Dictionary
\chapter*{How to Use This Dictionary}
\addcontentsline{toc}{chapter}{How to Use This Dictionary}

\section*{Entry Structure}

Each dictionary entry follows a consistent format:

\begin{itemize}[leftmargin=2cm]
    \item \textbf{Term Name} The main term being defined, in bold.
    \item \textit{Pronunciation and Part of Speech} Phonetic pronunciation followed by grammatical classification (noun, verb, abbreviation, etc.).
    \item \textbf{Definition} A clear, plain-language explanation of the term.
    \item \textit{Example} Real-world usage examples showing the term in context.
    \item \textit{See also} Cross-references to related terms within the dictionary.
    \item \textit{Etymology} Historical context or origin of the term (when relevant).
\end{itemize}

\section*{Numbered Definitions}

Some terms have multiple meanings. When a term has several distinct definitions, they are numbered:

\textbf{1.} First definition

\textbf{2.} Second definition

\section*{Cross-References}

Related terms are indicated with ``See also'' references. These cross-references help you explore connected concepts and build a comprehensive understanding of payment terminology.

\section*{Index}

Use the alphabetical index at the back of this dictionary to quickly locate any term. The index includes page numbers for easy navigation.

\section*{Abbreviations Used}

Common abbreviations in entries:
\begin{itemize}
    \item \textit{abbr.} = abbreviation
    \item \textit{n.} = noun
    \item \textit{v.} = verb
    \item \textit{adj.} = adjective
    \item \textit{pl.} = plural
\end{itemize}

For a complete list of abbreviations, see the Abbreviations section in the back matter.
